\documentclass[11pt,letterpaper]{article}
\usepackage[technical-reference]{cornell-notes}
\usepackage{needspace}

\title{Clause 17: Language Support Library - Cornell Notes}
\author{}
\date{}
\setDocTitle{Clause 17: Language Support Library}
\setDocSubtitle{C++ 2024 Cornell Notes}
\setDocOwner{C++ 2024 Cornell Notes}
\setDocAuthor{}
\setDocDate{}
\setCornellCollection{C++ 2024 Cornell Notes}
\setCornellUnitType{Clause}
\setCornellUnitNumber{17}
\setCornellUnitTitle{Language Support Library}
\setCornellCompanionLabel{C++ Technical Reference}
\hypersetup{pdftitle={Clause 17: Language Support Library - Cornell Notes},pdfsubject={C++ technical study notes},pdfauthor={}}

\begin{document}
\maketitle
\begin{CornellReferenceOrientationBox}
\textbf{Purpose.} Collects facilities closely tied to core execution: common types, implementation properties, numeric limits, startup and termination, dynamic allocation, type identification, source location, exceptions, initializer lists, comparisons, coroutines, and runtime support.
\par\textbf{Role.} The bridge between core-language semantics and essential library interfaces.
\par\textbf{Principal dependencies.} Uses the library-wide rules of Clause 16 together with the applicable core-language concepts and supporting library clauses.
\par\textbf{Primary audience.} programmers, runtime implementers, and low-level library authors.
\par\textbf{Major subject groups.} General; Common definitions; Implementation properties; Arithmetic types; Startup and termination; Dynamic memory management; Type identification; Source location; Exception handling; Initializer lists; Comparisons; Coroutines; and related subclauses.
\end{CornellReferenceOrientationBox}
\section{Learning Objectives}
\begin{itemize}[label=\textcolor{CornellReferenceTeal}{$\blacktriangleright$}]
\item Define the governing ideas of General and apply its stated contract at a call site.
\item Identify the governing ideas of Common definitions and apply its stated contract at a call site.
\item Distinguish the governing ideas of Implementation properties and apply its stated contract at a call site.
\item Explain the governing ideas of Arithmetic types and apply its stated contract at a call site.
\item Trace the governing ideas of Startup and termination and apply its stated contract at a call site.
\item Apply the governing ideas of Dynamic memory management and apply its stated contract at a call site.
\item Analyze the governing ideas of Type identification and apply its stated contract at a call site.
\item Evaluate the governing ideas of Source location and apply its stated contract at a call site.
\item Diagnose the governing ideas of Exception handling and apply its stated contract at a call site.
\item Compare the governing ideas of Initializer lists and apply its stated contract at a call site.
\item Verify the governing ideas of Comparisons and apply its stated contract at a call site.
\item Relate the governing ideas of Coroutines and apply its stated contract at a call site.
\item Classify the governing ideas of Other runtime support and apply its stated contract at a call site.
\item Justify the governing ideas of C headers and apply its stated contract at a call site.
\item Audit the governing ideas of Header <cstddef> synopsis and apply its stated contract at a call site.
\item Summarize the governing ideas of Header <cstdlib> synopsis and apply its stated contract at a call site.
\end{itemize}
\section{Normative Structure Map}
\small\begin{longtable}{@{}>{\RaggedRight\arraybackslash}p{0.13\textwidth}>{\RaggedRight\arraybackslash}p{0.27\textwidth}>{\RaggedRight\arraybackslash}p{0.19\textwidth}>{\RaggedRight\arraybackslash}p{0.31\textwidth}@{}}
\toprule\textbf{Subclause} & \textbf{Subject} & \textbf{Rule category} & \textbf{Key dependencies / entities}\\\midrule\endfirsthead
\toprule\textbf{Subclause} & \textbf{Subject} & \textbf{Rule category} & \textbf{Key dependencies / entities}\\\midrule\endhead
\bottomrule\endfoot
17.1 & General & library semantics & Uses the library-wide rules of Clause 16 together with the applicable core-language concepts and supporting library clauses.\\
17.2 & Common definitions & library semantics & Header <cstddef> synopsis, Header <cstdlib> synopsis, Null pointers, Sizes, alignments, and offsets, and related subclauses\\
17.3 & Implementation properties & library semantics & General, Header <version> synopsis, Header <limits> synopsis, Enum float\_\allowbreak{}round\_\allowbreak{}style, and related subclauses\\
17.4 & Arithmetic types & library semantics & Header <cstdint> synopsis, Header <stdfloat> synopsis\\
17.5 & Startup and termination & library semantics & Uses the library-wide rules of Clause 16 together with the applicable core-language concepts and supporting library clauses.\\
17.6 & Dynamic memory management & library semantics & General, Header <new> synopsis, Storage allocation and deallocation, Storage allocation errors, and related subclauses\\
17.7 & Type identification & library semantics & General, Header <typeinfo> synopsis, Class type\_\allowbreak{}info, Class bad\_\allowbreak{}cast, and related subclauses\\
17.8 & Source location & library semantics & Header <source\_\allowbreak{}location> synopsis, Class source\_\allowbreak{}location\\
17.9 & Exception handling & error behavior & General, Header <exception> synopsis, Class exception, Class bad\_\allowbreak{}exception, and related subclauses\\
17.10 & Initializer lists & library semantics & General, Header <initializer\_\allowbreak{}list> synopsis, Initializer list constructors, Initializer list access, and related subclauses\\
17.11 & Comparisons & library semantics & Header <compare> synopsis, Comparison category types, Class template common\_\allowbreak{}comparison\_\allowbreak{}category, Concept three\_\allowbreak{}way\_\allowbreak{}comparable, and related subclauses\\
17.12 & Coroutines & library semantics & General, Header <coroutine> synopsis, Coroutine traits, Class template coroutine\_\allowbreak{}handle, and related subclauses\\
17.13 & Other runtime support & library semantics & General, Header <cstdarg> synopsis, Header <csetjmp> synopsis, Header <csignal> synopsis, and related subclauses\\
17.14 & C headers & interface / declarations & General, Header <complex.h> synopsis, Header <iso646.h> synopsis, Header <stdalign.h> synopsis, and related subclauses\\
\end{longtable}\normalsize
\begin{CornellReferenceNormativeBox}A heading maps the clause; it does not replace the detailed conditions. For each operation, read requirements, effects, result, exceptions, complexity, remarks, and notes with their stated normative force.\end{CornellReferenceNormativeBox}
\subsection*{Rule Classification Guide}
\small\begin{longtable}{@{}>{\RaggedRight\arraybackslash}p{0.25\textwidth}>{\RaggedRight\arraybackslash}p{0.67\textwidth}@{}}\toprule\textbf{Label or category} & \textbf{How to read it}\\\midrule\endfirsthead\toprule\textbf{Label or category} & \textbf{How to read it (continued)}\\\midrule\endhead\bottomrule\endfoot
Constraints & Conditions controlling whether a declaration, overload, or specialization participates in the relevant context.\\
Mandates & Requirements checked for the described declaration or specialization; failure has the stated well-formedness consequence.\\
Preconditions & Obligations the caller establishes before evaluation; violating one forfeits the operation's contract.\\
Effects / Returns / Postconditions & Separate the state change, produced result, and state guaranteed after successful completion.\\
Throws / Error conditions & State how failure is reported and which exceptional paths are permitted or required.\\
Complexity & Bounds or exact counts for the operations named by the specification.\\
Remarks / Recommended practice & Remarks refine applicability or participation; recommended practice guides implementations without silently becoming a program requirement.\\
Exposition-only & A device for describing semantics, not a public name on which a program can depend.\\
\end{longtable}\normalsize
\section{Cornell Cue-and-Notes}
\Needspace{8\baselineskip}
\subsection*{17.1 General}
\begin{longtable}{@{}>{\RaggedRight\arraybackslash\columncolor{CornellReferenceCueGray}}p{0.28\textwidth}>{\RaggedRight\arraybackslash}p{0.64\textwidth}@{}}
\rowcolor{CornellReferenceNavy}\color{white}\textbf{Cue / recall prompt} & \color{white}\textbf{Notes}\\\endfirsthead
\rowcolor{CornellReferenceNavy}\color{white}\textbf{Cue / recall prompt} & \color{white}\textbf{Notes (continued)}\\\endhead
\arrayrulecolor{CornellReferenceLineGray}
\textbf{What contract governs General?}\\[-1mm]{\scriptsize\CornellClauseRef{17.1}} & Frames the terminology, applicability, and relationships needed to interpret General; detailed rules remain in the following subclauses.\\\hline
\end{longtable}
\begin{CornellReferenceSummaryBox}General supplies the library semantics portion of this clause. The key review move is to connect its local wording to the clause-wide model, then classify each condition by normative force before relying on it.\end{CornellReferenceSummaryBox}
\Needspace{8\baselineskip}
\subsection*{17.2 Common definitions}
\begin{longtable}{@{}>{\RaggedRight\arraybackslash\columncolor{CornellReferenceCueGray}}p{0.28\textwidth}>{\RaggedRight\arraybackslash}p{0.64\textwidth}@{}}
\rowcolor{CornellReferenceNavy}\color{white}\textbf{Cue / recall prompt} & \color{white}\textbf{Notes}\\\endfirsthead
\rowcolor{CornellReferenceNavy}\color{white}\textbf{Cue / recall prompt} & \color{white}\textbf{Notes (continued)}\\\endhead
\arrayrulecolor{CornellReferenceLineGray}
\textbf{What contract governs Header <cstddef> synopsis?}\\[-1mm]{\scriptsize\CornellClauseRef{17.2.1}} & Catalogs the declarations made available by Header <cstddef> synopsis. The synopsis is an interface map; the detailed specifications supply constraints and behavior.\\\hline
\textbf{What contract governs Header <cstdlib> synopsis?}\\[-1mm]{\scriptsize\CornellClauseRef{17.2.2}} & Catalogs the declarations made available by Header <cstdlib> synopsis. The synopsis is an interface map; the detailed specifications supply constraints and behavior.\\\hline
\textbf{What contract governs Null pointers?}\\[-1mm]{\scriptsize\CornellClauseRef{17.2.3}} & Specifies the facility family Null pointers. Read its declarations together with mandates, preconditions, effects, results, exceptions, complexity, and lifetime rules.\\\hline
\textbf{What contract governs Sizes, alignments, and offsets?}\\[-1mm]{\scriptsize\CornellClauseRef{17.2.4}} & Specifies the facility family Sizes, alignments, and offsets. Read its declarations together with mandates, preconditions, effects, results, exceptions, complexity, and lifetime rules.\\\hline
\textbf{What contract governs byte type operations?}\\[-1mm]{\scriptsize\CornellClauseRef{17.2.5}} & Defines the observable result of byte type operations together with applicable iterator-domain, ordering, complexity, invalidation, and exception conditions.\\\hline
\end{longtable}
\begin{CornellReferenceSummaryBox}Common definitions supplies the library semantics portion of this clause. The key review move is to connect its local wording to the clause-wide model, then classify each condition by normative force before relying on it.\end{CornellReferenceSummaryBox}
\Needspace{8\baselineskip}
\subsection*{17.3 Implementation properties}
\begin{longtable}{@{}>{\RaggedRight\arraybackslash\columncolor{CornellReferenceCueGray}}p{0.28\textwidth}>{\RaggedRight\arraybackslash}p{0.64\textwidth}@{}}
\rowcolor{CornellReferenceNavy}\color{white}\textbf{Cue / recall prompt} & \color{white}\textbf{Notes}\\\endfirsthead
\rowcolor{CornellReferenceNavy}\color{white}\textbf{Cue / recall prompt} & \color{white}\textbf{Notes (continued)}\\\endhead
\arrayrulecolor{CornellReferenceLineGray}
\textbf{What contract governs General?}\\[-1mm]{\scriptsize\CornellClauseRef{17.3.1}} & Frames the terminology, applicability, and relationships needed to interpret General; detailed rules remain in the following subclauses.\\\hline
\textbf{What contract governs Header <version> synopsis?}\\[-1mm]{\scriptsize\CornellClauseRef{17.3.2}} & Catalogs the declarations made available by Header <version> synopsis. The synopsis is an interface map; the detailed specifications supply constraints and behavior.\\\hline
\textbf{What contract governs Header <limits> synopsis?}\\[-1mm]{\scriptsize\CornellClauseRef{17.3.3}} & Catalogs the declarations made available by Header <limits> synopsis. The synopsis is an interface map; the detailed specifications supply constraints and behavior.\\\hline
\textbf{What contract governs Enum float\_\allowbreak{}round\_\allowbreak{}style?}\\[-1mm]{\scriptsize\CornellClauseRef{17.3.4}} & Specifies the facility family Enum float\_\allowbreak{}round\_\allowbreak{}style. Read its declarations together with mandates, preconditions, effects, results, exceptions, complexity, and lifetime rules.\\\hline
\textbf{What contract governs Class template numeric\_\allowbreak{}limits?}\\[-1mm]{\scriptsize\CornellClauseRef{17.3.5}} & Specifies the facility family Class template numeric\_\allowbreak{}limits. Read its declarations together with mandates, preconditions, effects, results, exceptions, complexity, and lifetime rules.\\\hline
\textbf{What contract governs Header <climits> synopsis?}\\[-1mm]{\scriptsize\CornellClauseRef{17.3.6}} & Catalogs the declarations made available by Header <climits> synopsis. The synopsis is an interface map; the detailed specifications supply constraints and behavior.\\\hline
\textbf{What contract governs Header <cfloat> synopsis?}\\[-1mm]{\scriptsize\CornellClauseRef{17.3.7}} & Catalogs the declarations made available by Header <cfloat> synopsis. The synopsis is an interface map; the detailed specifications supply constraints and behavior.\\\hline
\end{longtable}
\begin{CornellReferenceSummaryBox}Implementation properties supplies the library semantics portion of this clause. The key review move is to connect its local wording to the clause-wide model, then classify each condition by normative force before relying on it.\end{CornellReferenceSummaryBox}
\Needspace{5\baselineskip}\paragraph{Detailed coverage ledger.}
\begin{description}[style=nextline,leftmargin=2.8cm,font=\normalfont\bfseries\color{CornellReferenceTeal}]
\item[17.3.5.1] \textbf{General.} Frames the terminology, applicability, and relationships needed to interpret General; detailed rules remain in the following subclauses.
\item[17.3.5.2] \textbf{numeric\_\allowbreak{}limits members.} Specifies the facility family numeric\_\allowbreak{}limits members. Read its declarations together with mandates, preconditions, effects, results, exceptions, complexity, and lifetime rules.
\item[17.3.5.3] \textbf{numeric\_\allowbreak{}limits specializations.} Specifies the facility family numeric\_\allowbreak{}limits specializations. Read its declarations together with mandates, preconditions, effects, results, exceptions, complexity, and lifetime rules.
\end{description}
\Needspace{8\baselineskip}
\subsection*{17.4 Arithmetic types}
\begin{longtable}{@{}>{\RaggedRight\arraybackslash\columncolor{CornellReferenceCueGray}}p{0.28\textwidth}>{\RaggedRight\arraybackslash}p{0.64\textwidth}@{}}
\rowcolor{CornellReferenceNavy}\color{white}\textbf{Cue / recall prompt} & \color{white}\textbf{Notes}\\\endfirsthead
\rowcolor{CornellReferenceNavy}\color{white}\textbf{Cue / recall prompt} & \color{white}\textbf{Notes (continued)}\\\endhead
\arrayrulecolor{CornellReferenceLineGray}
\textbf{What contract governs Header <cstdint> synopsis?}\\[-1mm]{\scriptsize\CornellClauseRef{17.4.1}} & Catalogs the declarations made available by Header <cstdint> synopsis. The synopsis is an interface map; the detailed specifications supply constraints and behavior.\\\hline
\textbf{What contract governs Header <stdfloat> synopsis?}\\[-1mm]{\scriptsize\CornellClauseRef{17.4.2}} & Catalogs the declarations made available by Header <stdfloat> synopsis. The synopsis is an interface map; the detailed specifications supply constraints and behavior.\\\hline
\end{longtable}
\begin{CornellReferenceSummaryBox}Arithmetic types supplies the library semantics portion of this clause. The key review move is to connect its local wording to the clause-wide model, then classify each condition by normative force before relying on it.\end{CornellReferenceSummaryBox}
\Needspace{8\baselineskip}
\subsection*{17.5 Startup and termination}
\begin{longtable}{@{}>{\RaggedRight\arraybackslash\columncolor{CornellReferenceCueGray}}p{0.28\textwidth}>{\RaggedRight\arraybackslash}p{0.64\textwidth}@{}}
\rowcolor{CornellReferenceNavy}\color{white}\textbf{Cue / recall prompt} & \color{white}\textbf{Notes}\\\endfirsthead
\rowcolor{CornellReferenceNavy}\color{white}\textbf{Cue / recall prompt} & \color{white}\textbf{Notes (continued)}\\\endhead
\arrayrulecolor{CornellReferenceLineGray}
\textbf{What contract governs Startup and termination?}\\[-1mm]{\scriptsize\CornellClauseRef{17.5}} & Specifies the facility family Startup and termination. Read its declarations together with mandates, preconditions, effects, results, exceptions, complexity, and lifetime rules.\\\hline
\end{longtable}
\begin{CornellReferenceSummaryBox}Startup and termination supplies the library semantics portion of this clause. The key review move is to connect its local wording to the clause-wide model, then classify each condition by normative force before relying on it.\end{CornellReferenceSummaryBox}
\Needspace{8\baselineskip}
\subsection*{17.6 Dynamic memory management}
\begin{longtable}{@{}>{\RaggedRight\arraybackslash\columncolor{CornellReferenceCueGray}}p{0.28\textwidth}>{\RaggedRight\arraybackslash}p{0.64\textwidth}@{}}
\rowcolor{CornellReferenceNavy}\color{white}\textbf{Cue / recall prompt} & \color{white}\textbf{Notes}\\\endfirsthead
\rowcolor{CornellReferenceNavy}\color{white}\textbf{Cue / recall prompt} & \color{white}\textbf{Notes (continued)}\\\endhead
\arrayrulecolor{CornellReferenceLineGray}
\textbf{What contract governs General?}\\[-1mm]{\scriptsize\CornellClauseRef{17.6.1}} & Frames the terminology, applicability, and relationships needed to interpret General; detailed rules remain in the following subclauses.\\\hline
\textbf{What contract governs Header <new> synopsis?}\\[-1mm]{\scriptsize\CornellClauseRef{17.6.2}} & Catalogs the declarations made available by Header <new> synopsis. The synopsis is an interface map; the detailed specifications supply constraints and behavior.\\\hline
\textbf{What contract governs Storage allocation and deallocation?}\\[-1mm]{\scriptsize\CornellClauseRef{17.6.3}} & Specifies the facility family Storage allocation and deallocation. Read its declarations together with mandates, preconditions, effects, results, exceptions, complexity, and lifetime rules.\\\hline
\textbf{What contract governs Storage allocation errors?}\\[-1mm]{\scriptsize\CornellClauseRef{17.6.4}} & Defines the failure model for Storage allocation errors, including how an error is represented, reported, propagated, or converted into a diagnostic consequence.\\\hline
\textbf{What contract governs Pointer optimization barrier?}\\[-1mm]{\scriptsize\CornellClauseRef{17.6.5}} & Defines the blocking and synchronization contract for Pointer optimization barrier; ownership, wake conditions, timeout behavior, and happens-before edges must be analyzed together.\\\hline
\textbf{What contract governs Hardware interference size?}\\[-1mm]{\scriptsize\CornellClauseRef{17.6.6}} & Specifies the facility family Hardware interference size. Read its declarations together with mandates, preconditions, effects, results, exceptions, complexity, and lifetime rules.\\\hline
\end{longtable}
\begin{CornellReferenceSummaryBox}Dynamic memory management supplies the library semantics portion of this clause. The key review move is to connect its local wording to the clause-wide model, then classify each condition by normative force before relying on it.\end{CornellReferenceSummaryBox}
\Needspace{5\baselineskip}\paragraph{Detailed coverage ledger.}
\begin{description}[style=nextline,leftmargin=2.8cm,font=\normalfont\bfseries\color{CornellReferenceTeal}]
\item[17.6.3.1] \textbf{General.} Frames the terminology, applicability, and relationships needed to interpret General; detailed rules remain in the following subclauses.
\item[17.6.3.2] \textbf{Single-object forms.} Specifies the facility family Single-object forms. Read its declarations together with mandates, preconditions, effects, results, exceptions, complexity, and lifetime rules.
\item[17.6.3.3] \textbf{Array forms.} Specifies the facility family Array forms. Read its declarations together with mandates, preconditions, effects, results, exceptions, complexity, and lifetime rules.
\item[17.6.3.4] \textbf{Non-allocating forms.} Specifies the facility family Non-allocating forms. Read its declarations together with mandates, preconditions, effects, results, exceptions, complexity, and lifetime rules.
\item[17.6.3.5] \textbf{Data races.} Specifies the facility family Data races. Read its declarations together with mandates, preconditions, effects, results, exceptions, complexity, and lifetime rules.
\item[17.6.4.1] \textbf{Class bad\_\allowbreak{}alloc.} Specifies the facility family Class bad\_\allowbreak{}alloc. Read its declarations together with mandates, preconditions, effects, results, exceptions, complexity, and lifetime rules.
\item[17.6.4.2] \textbf{Class bad\_\allowbreak{}array\_\allowbreak{}new\_\allowbreak{}length.} Specifies the facility family Class bad\_\allowbreak{}array\_\allowbreak{}new\_\allowbreak{}length. Read its declarations together with mandates, preconditions, effects, results, exceptions, complexity, and lifetime rules.
\item[17.6.4.3] \textbf{Type new\_\allowbreak{}handler.} Specifies the facility family Type new\_\allowbreak{}handler. Read its declarations together with mandates, preconditions, effects, results, exceptions, complexity, and lifetime rules.
\item[17.6.4.4] \textbf{set\_\allowbreak{}new\_\allowbreak{}handler.} Specifies the facility family set\_\allowbreak{}new\_\allowbreak{}handler. Read its declarations together with mandates, preconditions, effects, results, exceptions, complexity, and lifetime rules.
\item[17.6.4.5] \textbf{get\_\allowbreak{}new\_\allowbreak{}handler.} Specifies the facility family get\_\allowbreak{}new\_\allowbreak{}handler. Read its declarations together with mandates, preconditions, effects, results, exceptions, complexity, and lifetime rules.
\end{description}
\Needspace{8\baselineskip}
\subsection*{17.7 Type identification}
\begin{longtable}{@{}>{\RaggedRight\arraybackslash\columncolor{CornellReferenceCueGray}}p{0.28\textwidth}>{\RaggedRight\arraybackslash}p{0.64\textwidth}@{}}
\rowcolor{CornellReferenceNavy}\color{white}\textbf{Cue / recall prompt} & \color{white}\textbf{Notes}\\\endfirsthead
\rowcolor{CornellReferenceNavy}\color{white}\textbf{Cue / recall prompt} & \color{white}\textbf{Notes (continued)}\\\endhead
\arrayrulecolor{CornellReferenceLineGray}
\textbf{What contract governs General?}\\[-1mm]{\scriptsize\CornellClauseRef{17.7.1}} & Frames the terminology, applicability, and relationships needed to interpret General; detailed rules remain in the following subclauses.\\\hline
\textbf{What contract governs Header <typeinfo> synopsis?}\\[-1mm]{\scriptsize\CornellClauseRef{17.7.2}} & Catalogs the declarations made available by Header <typeinfo> synopsis. The synopsis is an interface map; the detailed specifications supply constraints and behavior.\\\hline
\textbf{What contract governs Class type\_\allowbreak{}info?}\\[-1mm]{\scriptsize\CornellClauseRef{17.7.3}} & Specifies the facility family Class type\_\allowbreak{}info. Read its declarations together with mandates, preconditions, effects, results, exceptions, complexity, and lifetime rules.\\\hline
\textbf{What contract governs Class bad\_\allowbreak{}cast?}\\[-1mm]{\scriptsize\CornellClauseRef{17.7.4}} & Specifies the facility family Class bad\_\allowbreak{}cast. Read its declarations together with mandates, preconditions, effects, results, exceptions, complexity, and lifetime rules.\\\hline
\textbf{What contract governs Class bad\_\allowbreak{}typeid?}\\[-1mm]{\scriptsize\CornellClauseRef{17.7.5}} & Specifies the facility family Class bad\_\allowbreak{}typeid. Read its declarations together with mandates, preconditions, effects, results, exceptions, complexity, and lifetime rules.\\\hline
\end{longtable}
\begin{CornellReferenceSummaryBox}Type identification supplies the library semantics portion of this clause. The key review move is to connect its local wording to the clause-wide model, then classify each condition by normative force before relying on it.\end{CornellReferenceSummaryBox}
\Needspace{8\baselineskip}
\subsection*{17.8 Source location}
\begin{longtable}{@{}>{\RaggedRight\arraybackslash\columncolor{CornellReferenceCueGray}}p{0.28\textwidth}>{\RaggedRight\arraybackslash}p{0.64\textwidth}@{}}
\rowcolor{CornellReferenceNavy}\color{white}\textbf{Cue / recall prompt} & \color{white}\textbf{Notes}\\\endfirsthead
\rowcolor{CornellReferenceNavy}\color{white}\textbf{Cue / recall prompt} & \color{white}\textbf{Notes (continued)}\\\endhead
\arrayrulecolor{CornellReferenceLineGray}
\textbf{What contract governs Header <source\_\allowbreak{}location> synopsis?}\\[-1mm]{\scriptsize\CornellClauseRef{17.8.1}} & Catalogs the declarations made available by Header <source\_\allowbreak{}location> synopsis. The synopsis is an interface map; the detailed specifications supply constraints and behavior.\\\hline
\textbf{What contract governs Class source\_\allowbreak{}location?}\\[-1mm]{\scriptsize\CornellClauseRef{17.8.2}} & Specifies the facility family Class source\_\allowbreak{}location. Read its declarations together with mandates, preconditions, effects, results, exceptions, complexity, and lifetime rules.\\\hline
\end{longtable}
\begin{CornellReferenceSummaryBox}Source location supplies the library semantics portion of this clause. The key review move is to connect its local wording to the clause-wide model, then classify each condition by normative force before relying on it.\end{CornellReferenceSummaryBox}
\Needspace{5\baselineskip}\paragraph{Detailed coverage ledger.}
\begin{description}[style=nextline,leftmargin=2.8cm,font=\normalfont\bfseries\color{CornellReferenceTeal}]
\item[17.8.2.1] \textbf{General.} Frames the terminology, applicability, and relationships needed to interpret General; detailed rules remain in the following subclauses.
\item[17.8.2.2] \textbf{Creation.} Specifies the facility family Creation. Read its declarations together with mandates, preconditions, effects, results, exceptions, complexity, and lifetime rules.
\item[17.8.2.3] \textbf{Observers.} Specifies the facility family Observers. Read its declarations together with mandates, preconditions, effects, results, exceptions, complexity, and lifetime rules.
\end{description}
\Needspace{8\baselineskip}
\subsection*{17.9 Exception handling}
\begin{longtable}{@{}>{\RaggedRight\arraybackslash\columncolor{CornellReferenceCueGray}}p{0.28\textwidth}>{\RaggedRight\arraybackslash}p{0.64\textwidth}@{}}
\rowcolor{CornellReferenceNavy}\color{white}\textbf{Cue / recall prompt} & \color{white}\textbf{Notes}\\\endfirsthead
\rowcolor{CornellReferenceNavy}\color{white}\textbf{Cue / recall prompt} & \color{white}\textbf{Notes (continued)}\\\endhead
\arrayrulecolor{CornellReferenceLineGray}
\textbf{What contract governs General?}\\[-1mm]{\scriptsize\CornellClauseRef{17.9.1}} & Frames the terminology, applicability, and relationships needed to interpret General; detailed rules remain in the following subclauses.\\\hline
\textbf{What contract governs Header <exception> synopsis?}\\[-1mm]{\scriptsize\CornellClauseRef{17.9.2}} & Catalogs the declarations made available by Header <exception> synopsis. The synopsis is an interface map; the detailed specifications supply constraints and behavior.\\\hline
\textbf{What contract governs Class exception?}\\[-1mm]{\scriptsize\CornellClauseRef{17.9.3}} & Defines the failure model for Class exception, including how an error is represented, reported, propagated, or converted into a diagnostic consequence.\\\hline
\textbf{What contract governs Class bad\_\allowbreak{}exception?}\\[-1mm]{\scriptsize\CornellClauseRef{17.9.4}} & Defines the failure model for Class bad\_\allowbreak{}exception, including how an error is represented, reported, propagated, or converted into a diagnostic consequence.\\\hline
\textbf{What contract governs Abnormal termination?}\\[-1mm]{\scriptsize\CornellClauseRef{17.9.5}} & Specifies the facility family Abnormal termination. Read its declarations together with mandates, preconditions, effects, results, exceptions, complexity, and lifetime rules.\\\hline
\textbf{What contract governs uncaught\_\allowbreak{}exceptions?}\\[-1mm]{\scriptsize\CornellClauseRef{17.9.6}} & Defines the failure model for uncaught\_\allowbreak{}exceptions, including how an error is represented, reported, propagated, or converted into a diagnostic consequence.\\\hline
\textbf{What contract governs Exception propagation?}\\[-1mm]{\scriptsize\CornellClauseRef{17.9.7}} & Defines the failure model for Exception propagation, including how an error is represented, reported, propagated, or converted into a diagnostic consequence.\\\hline
\textbf{What contract governs nested\_\allowbreak{}exception?}\\[-1mm]{\scriptsize\CornellClauseRef{17.9.8}} & Defines the failure model for nested\_\allowbreak{}exception, including how an error is represented, reported, propagated, or converted into a diagnostic consequence.\\\hline
\end{longtable}
\begin{CornellReferenceSummaryBox}Exception handling supplies the error behavior portion of this clause. The key review move is to connect its local wording to the clause-wide model, then classify each condition by normative force before relying on it.\end{CornellReferenceSummaryBox}
\Needspace{5\baselineskip}\paragraph{Detailed coverage ledger.}
\begin{description}[style=nextline,leftmargin=2.8cm,font=\normalfont\bfseries\color{CornellReferenceTeal}]
\item[17.9.5.1] \textbf{Type terminate\_\allowbreak{}handler.} Specifies the facility family Type terminate\_\allowbreak{}handler. Read its declarations together with mandates, preconditions, effects, results, exceptions, complexity, and lifetime rules.
\item[17.9.5.2] \textbf{set\_\allowbreak{}terminate.} Specifies the facility family set\_\allowbreak{}terminate. Read its declarations together with mandates, preconditions, effects, results, exceptions, complexity, and lifetime rules.
\item[17.9.5.3] \textbf{get\_\allowbreak{}terminate.} Specifies the facility family get\_\allowbreak{}terminate. Read its declarations together with mandates, preconditions, effects, results, exceptions, complexity, and lifetime rules.
\item[17.9.5.4] \textbf{terminate.} Specifies the facility family terminate. Read its declarations together with mandates, preconditions, effects, results, exceptions, complexity, and lifetime rules.
\end{description}
\Needspace{8\baselineskip}
\subsection*{17.10 Initializer lists}
\begin{longtable}{@{}>{\RaggedRight\arraybackslash\columncolor{CornellReferenceCueGray}}p{0.28\textwidth}>{\RaggedRight\arraybackslash}p{0.64\textwidth}@{}}
\rowcolor{CornellReferenceNavy}\color{white}\textbf{Cue / recall prompt} & \color{white}\textbf{Notes}\\\endfirsthead
\rowcolor{CornellReferenceNavy}\color{white}\textbf{Cue / recall prompt} & \color{white}\textbf{Notes (continued)}\\\endhead
\arrayrulecolor{CornellReferenceLineGray}
\textbf{What contract governs General?}\\[-1mm]{\scriptsize\CornellClauseRef{17.10.1}} & Frames the terminology, applicability, and relationships needed to interpret General; detailed rules remain in the following subclauses.\\\hline
\textbf{What contract governs Header <initializer\_\allowbreak{}list> synopsis?}\\[-1mm]{\scriptsize\CornellClauseRef{17.10.2}} & Catalogs the declarations made available by Header <initializer\_\allowbreak{}list> synopsis. The synopsis is an interface map; the detailed specifications supply constraints and behavior.\\\hline
\textbf{How is Initializer list constructors initialized?}\\[-1mm]{\scriptsize\CornellClauseRef{17.10.3}} & Specifies how Initializer list constructors establishes object state, including participation constraints, initialization effects, and exceptional exits.\\\hline
\textbf{What contract governs Initializer list access?}\\[-1mm]{\scriptsize\CornellClauseRef{17.10.4}} & Determines the initialization form used by Initializer list access, the candidates or conversions considered, object-lifetime start, narrowing rules, and failure consequences.\\\hline
\textbf{What contract governs Initializer list range access?}\\[-1mm]{\scriptsize\CornellClauseRef{17.10.5}} & Explains the traversal or view contract for Initializer list range access; prove domain validity, lifetime, sentinel compatibility, and any borrowing or invalidation property.\\\hline
\end{longtable}
\begin{CornellReferenceSummaryBox}Initializer lists supplies the library semantics portion of this clause. The key review move is to connect its local wording to the clause-wide model, then classify each condition by normative force before relying on it.\end{CornellReferenceSummaryBox}
\Needspace{8\baselineskip}
\subsection*{17.11 Comparisons}
\begin{longtable}{@{}>{\RaggedRight\arraybackslash\columncolor{CornellReferenceCueGray}}p{0.28\textwidth}>{\RaggedRight\arraybackslash}p{0.64\textwidth}@{}}
\rowcolor{CornellReferenceNavy}\color{white}\textbf{Cue / recall prompt} & \color{white}\textbf{Notes}\\\endfirsthead
\rowcolor{CornellReferenceNavy}\color{white}\textbf{Cue / recall prompt} & \color{white}\textbf{Notes (continued)}\\\endhead
\arrayrulecolor{CornellReferenceLineGray}
\textbf{What contract governs Header <compare> synopsis?}\\[-1mm]{\scriptsize\CornellClauseRef{17.11.1}} & Catalogs the declarations made available by Header <compare> synopsis. The synopsis is an interface map; the detailed specifications supply constraints and behavior.\\\hline
\textbf{What contract governs Comparison category types?}\\[-1mm]{\scriptsize\CornellClauseRef{17.11.2}} & Specifies the facility family Comparison category types. Read its declarations together with mandates, preconditions, effects, results, exceptions, complexity, and lifetime rules.\\\hline
\textbf{What contract governs Class template common\_\allowbreak{}comparison\_\allowbreak{}category?}\\[-1mm]{\scriptsize\CornellClauseRef{17.11.3}} & Specifies the facility family Class template common\_\allowbreak{}comparison\_\allowbreak{}category. Read its declarations together with mandates, preconditions, effects, results, exceptions, complexity, and lifetime rules.\\\hline
\textbf{What contract governs Concept three\_\allowbreak{}way\_\allowbreak{}comparable?}\\[-1mm]{\scriptsize\CornellClauseRef{17.11.4}} & Defines or applies the generic requirement Concept three\_\allowbreak{}way\_\allowbreak{}comparable; syntactic satisfaction must be paired with every stated semantic and equality-preservation expectation.\\\hline
\textbf{What contract governs Result of three-way comparison?}\\[-1mm]{\scriptsize\CornellClauseRef{17.11.5}} & Specifies the facility family Result of three-way comparison. Read its declarations together with mandates, preconditions, effects, results, exceptions, complexity, and lifetime rules.\\\hline
\textbf{What does Comparison algorithms guarantee?}\\[-1mm]{\scriptsize\CornellClauseRef{17.11.6}} & Defines the observable result of Comparison algorithms together with applicable iterator-domain, ordering, complexity, invalidation, and exception conditions.\\\hline
\end{longtable}
\begin{CornellReferenceSummaryBox}Comparisons supplies the library semantics portion of this clause. The key review move is to connect its local wording to the clause-wide model, then classify each condition by normative force before relying on it.\end{CornellReferenceSummaryBox}
\Needspace{5\baselineskip}\paragraph{Detailed coverage ledger.}
\begin{description}[style=nextline,leftmargin=2.8cm,font=\normalfont\bfseries\color{CornellReferenceTeal}]
\item[17.11.2.1] \textbf{Preamble.} Frames the terminology, applicability, and relationships needed to interpret Preamble; detailed rules remain in the following subclauses.
\item[17.11.2.2] \textbf{Class partial\_\allowbreak{}ordering.} Specifies the facility family Class partial\_\allowbreak{}ordering. Read its declarations together with mandates, preconditions, effects, results, exceptions, complexity, and lifetime rules.
\item[17.11.2.3] \textbf{Class weak\_\allowbreak{}ordering.} Specifies the facility family Class weak\_\allowbreak{}ordering. Read its declarations together with mandates, preconditions, effects, results, exceptions, complexity, and lifetime rules.
\item[17.11.2.4] \textbf{Class strong\_\allowbreak{}ordering.} Specifies the facility family Class strong\_\allowbreak{}ordering. Read its declarations together with mandates, preconditions, effects, results, exceptions, complexity, and lifetime rules.
\end{description}
\Needspace{8\baselineskip}
\subsection*{17.12 Coroutines}
\begin{longtable}{@{}>{\RaggedRight\arraybackslash\columncolor{CornellReferenceCueGray}}p{0.28\textwidth}>{\RaggedRight\arraybackslash}p{0.64\textwidth}@{}}
\rowcolor{CornellReferenceNavy}\color{white}\textbf{Cue / recall prompt} & \color{white}\textbf{Notes}\\\endfirsthead
\rowcolor{CornellReferenceNavy}\color{white}\textbf{Cue / recall prompt} & \color{white}\textbf{Notes (continued)}\\\endhead
\arrayrulecolor{CornellReferenceLineGray}
\textbf{What contract governs General?}\\[-1mm]{\scriptsize\CornellClauseRef{17.12.1}} & Frames the terminology, applicability, and relationships needed to interpret General; detailed rules remain in the following subclauses.\\\hline
\textbf{What contract governs Header <coroutine> synopsis?}\\[-1mm]{\scriptsize\CornellClauseRef{17.12.2}} & Catalogs the declarations made available by Header <coroutine> synopsis. The synopsis is an interface map; the detailed specifications supply constraints and behavior.\\\hline
\textbf{What contract governs Coroutine traits?}\\[-1mm]{\scriptsize\CornellClauseRef{17.12.3}} & Specifies the facility family Coroutine traits. Read its declarations together with mandates, preconditions, effects, results, exceptions, complexity, and lifetime rules.\\\hline
\textbf{What contract governs Class template coroutine\_\allowbreak{}handle?}\\[-1mm]{\scriptsize\CornellClauseRef{17.12.4}} & Specifies the facility family Class template coroutine\_\allowbreak{}handle. Read its declarations together with mandates, preconditions, effects, results, exceptions, complexity, and lifetime rules.\\\hline
\textbf{What contract governs No-op coroutines?}\\[-1mm]{\scriptsize\CornellClauseRef{17.12.5}} & Specifies the facility family No-op coroutines. Read its declarations together with mandates, preconditions, effects, results, exceptions, complexity, and lifetime rules.\\\hline
\textbf{What contract governs Trivial awaitables?}\\[-1mm]{\scriptsize\CornellClauseRef{17.12.6}} & Specifies the facility family Trivial awaitables. Read its declarations together with mandates, preconditions, effects, results, exceptions, complexity, and lifetime rules.\\\hline
\end{longtable}
\begin{CornellReferenceSummaryBox}Coroutines supplies the library semantics portion of this clause. The key review move is to connect its local wording to the clause-wide model, then classify each condition by normative force before relying on it.\end{CornellReferenceSummaryBox}
\Needspace{5\baselineskip}\paragraph{Detailed coverage ledger.}
\begin{description}[style=nextline,leftmargin=2.8cm,font=\normalfont\bfseries\color{CornellReferenceTeal}]
\item[17.12.3.1] \textbf{General.} Frames the terminology, applicability, and relationships needed to interpret General; detailed rules remain in the following subclauses.
\item[17.12.3.2] \textbf{Class template coroutine\_\allowbreak{}traits.} Specifies the facility family Class template coroutine\_\allowbreak{}traits. Read its declarations together with mandates, preconditions, effects, results, exceptions, complexity, and lifetime rules.
\item[17.12.4.1] \textbf{General.} Frames the terminology, applicability, and relationships needed to interpret General; detailed rules remain in the following subclauses.
\item[17.12.4.2] \textbf{Construct/reset.} Specifies the facility family Construct/reset. Read its declarations together with mandates, preconditions, effects, results, exceptions, complexity, and lifetime rules.
\item[17.12.4.3] \textbf{Conversion.} Specifies source and destination conditions for Conversion, the resulting type or value category, and any information-loss, pointer, qualification, or lifetime consequence.
\item[17.12.4.4] \textbf{Export/import.} Specifies the facility family Export/import. Read its declarations together with mandates, preconditions, effects, results, exceptions, complexity, and lifetime rules.
\item[17.12.4.5] \textbf{Observers.} Specifies the facility family Observers. Read its declarations together with mandates, preconditions, effects, results, exceptions, complexity, and lifetime rules.
\item[17.12.4.6] \textbf{Resumption.} Specifies the facility family Resumption. Read its declarations together with mandates, preconditions, effects, results, exceptions, complexity, and lifetime rules.
\item[17.12.4.7] \textbf{Promise access.} Defines asynchronous shared-state behavior for Promise access; establish producer/consumer ownership, readiness, blocking, exception transport, and destruction effects.
\item[17.12.4.8] \textbf{Comparison operators.} Specifies the facility family Comparison operators. Read its declarations together with mandates, preconditions, effects, results, exceptions, complexity, and lifetime rules.
\item[17.12.4.9] \textbf{Hash support.} Specifies the facility family Hash support. Read its declarations together with mandates, preconditions, effects, results, exceptions, complexity, and lifetime rules.
\item[17.12.5.1] \textbf{Class noop\_\allowbreak{}coroutine\_\allowbreak{}promise.} Defines asynchronous shared-state behavior for Class noop\_\allowbreak{}coroutine\_\allowbreak{}promise; establish producer/consumer ownership, readiness, blocking, exception transport, and destruction effects.
\item[17.12.5.2] \textbf{Class coroutine\_\allowbreak{}handle<noop\_\allowbreak{}coroutine\_\allowbreak{}promise>.} Defines asynchronous shared-state behavior for Class coroutine\_\allowbreak{}handle<noop\_\allowbreak{}coroutine\_\allowbreak{}promise>; establish producer/consumer ownership, readiness, blocking, exception transport, and destruction effects.
\item[17.12.5.2.1] \textbf{General.} Frames the terminology, applicability, and relationships needed to interpret General; detailed rules remain in the following subclauses.
\item[17.12.5.2.2] \textbf{Conversion.} Specifies source and destination conditions for Conversion, the resulting type or value category, and any information-loss, pointer, qualification, or lifetime consequence.
\item[17.12.5.2.3] \textbf{Observers.} Specifies the facility family Observers. Read its declarations together with mandates, preconditions, effects, results, exceptions, complexity, and lifetime rules.
\item[17.12.5.2.4] \textbf{Resumption.} Specifies the facility family Resumption. Read its declarations together with mandates, preconditions, effects, results, exceptions, complexity, and lifetime rules.
\item[17.12.5.2.5] \textbf{Promise access.} Defines asynchronous shared-state behavior for Promise access; establish producer/consumer ownership, readiness, blocking, exception transport, and destruction effects.
\item[17.12.5.2.6] \textbf{Address.} Specifies the facility family Address. Read its declarations together with mandates, preconditions, effects, results, exceptions, complexity, and lifetime rules.
\item[17.12.5.3] \textbf{Function noop\_\allowbreak{}coroutine.} Specifies the facility family Function noop\_\allowbreak{}coroutine. Read its declarations together with mandates, preconditions, effects, results, exceptions, complexity, and lifetime rules.
\end{description}
\Needspace{8\baselineskip}
\subsection*{17.13 Other runtime support}
\begin{longtable}{@{}>{\RaggedRight\arraybackslash\columncolor{CornellReferenceCueGray}}p{0.28\textwidth}>{\RaggedRight\arraybackslash}p{0.64\textwidth}@{}}
\rowcolor{CornellReferenceNavy}\color{white}\textbf{Cue / recall prompt} & \color{white}\textbf{Notes}\\\endfirsthead
\rowcolor{CornellReferenceNavy}\color{white}\textbf{Cue / recall prompt} & \color{white}\textbf{Notes (continued)}\\\endhead
\arrayrulecolor{CornellReferenceLineGray}
\textbf{What contract governs General?}\\[-1mm]{\scriptsize\CornellClauseRef{17.13.1}} & Frames the terminology, applicability, and relationships needed to interpret General; detailed rules remain in the following subclauses.\\\hline
\textbf{What contract governs Header <cstdarg> synopsis?}\\[-1mm]{\scriptsize\CornellClauseRef{17.13.2}} & Catalogs the declarations made available by Header <cstdarg> synopsis. The synopsis is an interface map; the detailed specifications supply constraints and behavior.\\\hline
\textbf{What contract governs Header <csetjmp> synopsis?}\\[-1mm]{\scriptsize\CornellClauseRef{17.13.3}} & Catalogs the declarations made available by Header <csetjmp> synopsis. The synopsis is an interface map; the detailed specifications supply constraints and behavior.\\\hline
\textbf{What contract governs Header <csignal> synopsis?}\\[-1mm]{\scriptsize\CornellClauseRef{17.13.4}} & Catalogs the declarations made available by Header <csignal> synopsis. The synopsis is an interface map; the detailed specifications supply constraints and behavior.\\\hline
\textbf{What contract governs Signal handlers?}\\[-1mm]{\scriptsize\CornellClauseRef{17.13.5}} & Specifies the facility family Signal handlers. Read its declarations together with mandates, preconditions, effects, results, exceptions, complexity, and lifetime rules.\\\hline
\end{longtable}
\begin{CornellReferenceSummaryBox}Other runtime support supplies the library semantics portion of this clause. The key review move is to connect its local wording to the clause-wide model, then classify each condition by normative force before relying on it.\end{CornellReferenceSummaryBox}
\Needspace{8\baselineskip}
\subsection*{17.14 C headers}
\begin{longtable}{@{}>{\RaggedRight\arraybackslash\columncolor{CornellReferenceCueGray}}p{0.28\textwidth}>{\RaggedRight\arraybackslash}p{0.64\textwidth}@{}}
\rowcolor{CornellReferenceNavy}\color{white}\textbf{Cue / recall prompt} & \color{white}\textbf{Notes}\\\endfirsthead
\rowcolor{CornellReferenceNavy}\color{white}\textbf{Cue / recall prompt} & \color{white}\textbf{Notes (continued)}\\\endhead
\arrayrulecolor{CornellReferenceLineGray}
\textbf{What contract governs General?}\\[-1mm]{\scriptsize\CornellClauseRef{17.14.1}} & Frames the terminology, applicability, and relationships needed to interpret General; detailed rules remain in the following subclauses.\\\hline
\textbf{What contract governs Header <complex.h> synopsis?}\\[-1mm]{\scriptsize\CornellClauseRef{17.14.2}} & Catalogs the declarations made available by Header <complex.h> synopsis. The synopsis is an interface map; the detailed specifications supply constraints and behavior.\\\hline
\textbf{What contract governs Header <iso646.h> synopsis?}\\[-1mm]{\scriptsize\CornellClauseRef{17.14.3}} & Catalogs the declarations made available by Header <iso646.h> synopsis. The synopsis is an interface map; the detailed specifications supply constraints and behavior.\\\hline
\textbf{What contract governs Header <stdalign.h> synopsis?}\\[-1mm]{\scriptsize\CornellClauseRef{17.14.4}} & Catalogs the declarations made available by Header <stdalign.h> synopsis. The synopsis is an interface map; the detailed specifications supply constraints and behavior.\\\hline
\textbf{What contract governs Header <stdbool.h> synopsis?}\\[-1mm]{\scriptsize\CornellClauseRef{17.14.5}} & Catalogs the declarations made available by Header <stdbool.h> synopsis. The synopsis is an interface map; the detailed specifications supply constraints and behavior.\\\hline
\textbf{What contract governs Header <tgmath.h> synopsis?}\\[-1mm]{\scriptsize\CornellClauseRef{17.14.6}} & Catalogs the declarations made available by Header <tgmath.h> synopsis. The synopsis is an interface map; the detailed specifications supply constraints and behavior.\\\hline
\textbf{What contract governs Other C headers?}\\[-1mm]{\scriptsize\CornellClauseRef{17.14.7}} & Specifies the facility family Other C headers. Read its declarations together with mandates, preconditions, effects, results, exceptions, complexity, and lifetime rules.\\\hline
\end{longtable}
\begin{CornellReferenceSummaryBox}C headers supplies the interface / declarations portion of this clause. The key review move is to connect its local wording to the clause-wide model, then classify each condition by normative force before relying on it.\end{CornellReferenceSummaryBox}
\section{Language-Lawyer and Library-Usage Pitfalls}
\begin{CornellReferencePitfallBox}\begin{itemize}[label=\textcolor{CornellReferenceAmber}{$\blacktriangleright$}]
\item Reading a synopsis as the complete behavioral contract.
\item Ignoring mandates, preconditions, exception behavior, or complexity requirements.
\item Using an exposition-only name as though it were public API.
\item Assuming invalidation, ownership, or thread-safety properties that are not stated.
\item Treating successful compilation as proof that semantic requirements and preconditions are met.
\end{itemize}\end{CornellReferencePitfallBox}
\section{Programmer and Implementation Checklist}
\begin{CornellReferenceChecklistBox}\begin{itemize}[label=$\square$]
\item Identify the declaring header or module exposure for each facility.
\item Check template arguments and mandates before evaluating an operation.
\item Establish every precondition at the call site.
\item Record effects, returned value, postconditions, and exceptions separately.
\item Audit iterator, reference, pointer, and object lifetime after mutation.
\item Verify stated complexity and synchronization assumptions.
\item For \CornellClauseRef{17.1}, verify general against the precise rule category and all cited dependencies.
\item For \CornellClauseRef{17.2}, verify common definitions against the precise rule category and all cited dependencies.
\item For \CornellClauseRef{17.3}, verify implementation properties against the precise rule category and all cited dependencies.
\item For \CornellClauseRef{17.4}, verify arithmetic types against the precise rule category and all cited dependencies.
\end{itemize}\end{CornellReferenceChecklistBox}
\section{Clause Synthesis}
\begin{CornellReferenceOrientationBox}Collects facilities closely tied to core execution: common types, implementation properties, numeric limits, startup and termination, dynamic allocation, type identification, source location, exceptions, initializer lists, comparisons, coroutines, and runtime support. The bridge between core-language semantics and essential library interfaces. A sound review distinguishes syntax and participation from runtime preconditions, keeps notes and examples non-mandatory, and follows cross-clause dependencies before making a portability claim.\end{CornellReferenceOrientationBox}
\section{Self-Test Questions}
\begin{enumerate}
\item What role does General play in the clause's overall model?
\item Which normative distinctions must be kept separate when applying Common definitions?
\item How would you audit a program or implementation that depends on Implementation properties?
\item Compare Arithmetic types with Startup and termination; what different question does each answer?
\item What role does Startup and termination play in the clause's overall model?
\item Which normative distinctions must be kept separate when applying Dynamic memory management?
\item How would you audit a program or implementation that depends on Type identification?
\item Compare Source location with Exception handling; what different question does each answer?
\item What role does Exception handling play in the clause's overall model?
\item Which normative distinctions must be kept separate when applying Initializer lists?
\item How would you audit a program or implementation that depends on Comparisons?
\item Compare Coroutines with Other runtime support; what different question does each answer?
\item What role does Other runtime support play in the clause's overall model?
\item Which normative distinctions must be kept separate when applying C headers?
\item Scenario: a program relies on an unstated implementation behavior while using Header <cstddef> synopsis. What must be checked before calling the program portable?
\item Scenario: a program relies on an unstated implementation behavior while using Header <cstdlib> synopsis. What must be checked before calling the program portable?
\end{enumerate}
\clearpage\section{Answer Key}
\begin{enumerate}
\item General is one part of the clause's library semantics model. Its role is understood by connecting the local rule in 17.14.1 to the clause orientation and its dependent concepts.
\item Keep formation and well-formedness separate from runtime preconditions and effects. Also distinguish mandatory wording from notes, implementation choice from undefined behavior, and interface declarations from detailed contracts; see 17.2.
\item Audit Implementation properties by locating the controlling wording in 17.3, listing inputs and type constraints, proving any preconditions, recording effects and failure behavior, and checking lifetime, invalidation, complexity, and synchronization where applicable.
\item Arithmetic types and Startup and termination occupy different steps or facility groups. Analyze each under its own subclause, then follow the explicit dependency between them rather than transferring one set of guarantees to the other; begin at 17.4.
\item Startup and termination is one part of the clause's library semantics model. Its role is understood by connecting the local rule in 17.5 to the clause orientation and its dependent concepts.
\item Keep formation and well-formedness separate from runtime preconditions and effects. Also distinguish mandatory wording from notes, implementation choice from undefined behavior, and interface declarations from detailed contracts; see 17.6.
\item Audit Type identification by locating the controlling wording in 17.7, listing inputs and type constraints, proving any preconditions, recording effects and failure behavior, and checking lifetime, invalidation, complexity, and synchronization where applicable.
\item Source location and Exception handling occupy different steps or facility groups. Analyze each under its own subclause, then follow the explicit dependency between them rather than transferring one set of guarantees to the other; begin at 17.8.
\item Exception handling is one part of the clause's error behavior model. Its role is understood by connecting the local rule in 17.9 to the clause orientation and its dependent concepts.
\item Keep formation and well-formedness separate from runtime preconditions and effects. Also distinguish mandatory wording from notes, implementation choice from undefined behavior, and interface declarations from detailed contracts; see 17.10.
\item Audit Comparisons by locating the controlling wording in 17.11, listing inputs and type constraints, proving any preconditions, recording effects and failure behavior, and checking lifetime, invalidation, complexity, and synchronization where applicable.
\item Coroutines and Other runtime support occupy different steps or facility groups. Analyze each under its own subclause, then follow the explicit dependency between them rather than transferring one set of guarantees to the other; begin at 17.12.
\item Other runtime support is one part of the clause's library semantics model. Its role is understood by connecting the local rule in 17.13 to the clause orientation and its dependent concepts.
\item Keep formation and well-formedness separate from runtime preconditions and effects. Also distinguish mandatory wording from notes, implementation choice from undefined behavior, and interface declarations from detailed contracts; see 17.14.
\item First identify the exact requirement in 17.2.1, then classify the relied-on behavior as required, unspecified, implementation-defined, conditionally supported, or undefined. Portability requires satisfying the program obligations and avoiding dependence on a choice not guaranteed in the target environments.
\item First identify the exact requirement in 17.2.2, then classify the relied-on behavior as required, unspecified, implementation-defined, conditionally supported, or undefined. Portability requires satisfying the program obligations and avoiding dependence on a choice not guaranteed in the target environments.
\end{enumerate}
\section{Key-Term Recap}
\begin{longtable}{@{}>{\RaggedRight\arraybackslash}p{0.28\textwidth}>{\RaggedRight\arraybackslash}p{0.64\textwidth}@{}}\toprule\textbf{Term} & \textbf{Working meaning}\\\midrule\endfirsthead\toprule\textbf{Term} & \textbf{Working meaning (continued)}\\\midrule\endhead\bottomrule\endfoot
numeric\_\allowbreak{}limits & A type-traits interface describing arithmetic properties.\\
new\_\allowbreak{}handler & A program-installable handler used by allocation-failure processing.\\
source\_\allowbreak{}location & A value type representing a call-site source position.\\
exception\_\allowbreak{}ptr & A shareable handle to an exception object for later propagation.\\
strong ordering & A comparison category expressing substitutable equality and total ordering.\\
coroutine\_\allowbreak{}handle & A non-owning interface to a suspended coroutine state.\\
\end{longtable}
\CornellTakeaway{Language-support facilities expose implementation properties and runtime mechanisms that must remain consistent with core-language semantics.}
\end{document}
